\section{Chapter 10: Modules over a ring}\label{sec:14-appendix-solutions:09-chapter-10:1}

\textbf{1. The identity function is a linear map}

\begin{lstlisting}
def idLinearMap {R : Type} (Rg : Ring R) {M : Type} (Mod : Module R Rg M) :
    LinearMap Rg Mod Mod where
  toFun := id
  map_add := by
    intro m n
    rfl
  map_smul := by
    intro r m
    rfl
\end{lstlisting}

\texttt{id : M → M} is \texttt{fun x => x}. Both fields reduce to \href{https://lean-lang.org/doc/reference/latest/Tactic-Proofs/Tactic-Reference/}{\texttt{rfl}} because \texttt{id (Mod.addGrp.op m n)} and \texttt{Mod.addGrp.op (id m) (id n)} unfold to the same term (\texttt{id} changes nothing). The same holds for \texttt{map\_\allowbreak{}smul}. The only content in \texttt{LinearMap}'s two fields is the compatibility of \texttt{toFun} with \texttt{+}/\texttt{smul}, and \texttt{id} trivially preserves everything, since it changes nothing.

\textbf{2. Linear maps compose}

\begin{lstlisting}
def composeLinearMap {R : Type} (Rg : Ring R) {M N P : Type}
    {ModM : Module R Rg M} {ModN : Module R Rg N} {ModP : Module R Rg P}
    (f : LinearMap Rg ModM ModN) (g : LinearMap Rg ModN ModP) :
    LinearMap Rg ModM ModP where
  toFun := g.toFun ∘ f.toFun
  map_add := by
    intro m n
    show g.toFun (f.toFun (ModM.addGrp.op m n)) =
      ModP.addGrp.op (g.toFun (f.toFun m)) (g.toFun (f.toFun n))
    rw [f.map_add]
    -- Goal: g.toFun (op (f.toFun m) (f.toFun n)) = op (g.toFun (f.toFun m)) (g.toFun (f.toFun n))
    exact g.map_add (f.toFun m) (f.toFun n)
  map_smul := by
    intro r m
    show g.toFun (f.toFun (ModM.smul r m)) = ModP.smul r (g.toFun (f.toFun m))
    rw [f.map_smul]
    -- Goal: g.toFun (ModN.smul r (f.toFun m)) = ModP.smul r (g.toFun (f.toFun m))
    exact g.map_smul r (f.toFun m)
\end{lstlisting}

\texttt{g.toFun ∘ f.toFun} is ordinary function composition. Each field's proof first uses \texttt{f}'s own compatibility fact (\texttt{f.map\_\allowbreak{}add}/\texttt{f.map\_\allowbreak{}smul}) to push the additive-group operation or scalar action \emph{through} \texttt{f}. This leaves a goal that is now exactly \texttt{g}'s compatibility fact, applied to the already-transformed points \texttt{f.toFun m}, \texttt{f.toFun n}. This, plus Exercise 1 and the associativity of \texttt{∘} (both free from ordinary function composition), is exactly the data that makes \(R\)-modules and \(R\)-linear maps a category, as promised in the chapter text.

\textbf{3. Verifying \texttt{intSmul} satisfies \texttt{Module}'s axioms (partial: \texttt{one\_\allowbreak{}smul}, \texttt{smul\_\allowbreak{}add})}

\begin{lstlisting}
-- one_smul: intSmul CG 1 m = m
theorem intSmul_one_smul {M : Type} (CG : CommGroup M) (m : M) :
    intSmul CG 1 m = m := by
  show natSmul CG.toGroup 1 m = m
  show CG.toGroup.op m (natSmul CG.toGroup 0 m) = m
  show CG.toGroup.op m CG.toGroup.id = m
  exact CG.toGroup.id_right m
\end{lstlisting}

\texttt{intSmul CG 1 m} unfolds (since \texttt{1 = Int.ofNat 1}) to \texttt{natSmul CG.toGroup 1 m}, which unfolds (since \texttt{1 = Nat.succ Nat.zero}) to \texttt{CG.toGroup.op m (natSmul CG.toGroup 0 m)}. This unfolds again (the \texttt{zero} case of \texttt{natSmul}) to \texttt{CG.toGroup.op m CG.toGroup.id}, which is exactly the \texttt{id\_\allowbreak{}right} axiom.

\begin{lstlisting}
-- smul_add, restricted to natural-number scalars, by induction on n
theorem natSmul_add {M : Type} (Grp : Group M) (n : Nat) (m1 m2 : M)
    (comm : ∀ a b : M, Grp.op a b = Grp.op b a) :
    natSmul Grp n (Grp.op m1 m2) = Grp.op (natSmul Grp n m1) (natSmul Grp n m2) := by
  induction n with
  | zero =>
    -- Goal: natSmul Grp 0 (op m1 m2) = op (natSmul Grp 0 m1) (natSmul Grp 0 m2)
    show Grp.id = Grp.op Grp.id Grp.id
    exact (Grp.id_left Grp.id).symm
  | succ k ih =>
    -- ih : natSmul Grp k (op m1 m2) = op (natSmul Grp k m1) (natSmul Grp k m2)
    show Grp.op (Grp.op m1 m2) (natSmul Grp k (Grp.op m1 m2)) =
      Grp.op (Grp.op m1 (natSmul Grp k m1)) (Grp.op m2 (natSmul Grp k m2))
    rw [ih]
    -- Goal: op (op m1 m2) (op (natSmul Grp k m1) (natSmul Grp k m2))
    --     = op (op m1 (natSmul Grp k m1)) (op m2 (natSmul Grp k m2))
    -- Both sides are the same four elements combined via `op`, just grouped
    -- and ordered differently; regroup with `assoc` and swap the middle two
    -- terms using `comm`, exactly the "regroup, then rearrange" pattern
    -- from Chapter 7.
    rw [← Grp.assoc, Grp.assoc m1 m2, comm m2 (natSmul Grp k m1), ← Grp.assoc m1,
        Grp.assoc]
\end{lstlisting}

The base case is \texttt{id = id + id} (an identity is its own double). This mirrors the ``element equal to its own double is zero'' fact used in Chapter 9's \texttt{mul\_\allowbreak{}zero}. The inductive step requires commutativity (\texttt{comm}, available since the group is a \texttt{CommGroup}) to slide \texttt{m2} and \texttt{natSmul Grp k m1} past each other. This is exactly why \texttt{intSmul} distributing over \texttt{+} requires an abelian group, rather than a general one: without \texttt{comm}, there is no way to reorder the four terms into a matching shape. Full verification of \texttt{add\_\allowbreak{}smul} and \texttt{smul\_\allowbreak{}smul} follows the same induction-on-the-scalar pattern, and makes a good longer exercise for further practice.

\textbf{4. Submodule of multiples of \texttt{d}}

\begin{lstlisting}
def multiplesSubmodule (d : Int) : Submodule intRing intZModule where
  carrier := fun m => ∃ k : Int, m = d * k
  zero_mem := ⟨0, by show (0 : Int) = d * 0; rw [Int.mul_zero]⟩
  add_mem := by
    intro m n ⟨k, hk⟩ ⟨j, hj⟩
    refine ⟨k + j, ?_⟩
    show m + n = d * (k + j)
    rw [hk, hj, Int.mul_add]
  smul_mem := by
    intro r m ⟨k, hk⟩
    refine ⟨r * k, ?_⟩
    show r * m = d * (r * k)
    rw [hk, ← Int.mul_assoc, Int.mul_comm r d, Int.mul_assoc]
\end{lstlisting}

This has the same shape as \texttt{evenSubmodule} (the case \texttt{d = 2}). Every \texttt{2} in that proof is simply replaced by the parameter \texttt{d}, and each closure proof still reduces to an \texttt{Int} equation, handled the same way (\href{https://lean-lang.org/doc/reference/latest/Tactic-Proofs/Tactic-Reference/}{\texttt{show}} to reveal the goal's \texttt{+}/\texttt{*}-form, then \href{https://lean-lang.org/doc/reference/latest/Tactic-Proofs/Tactic-Reference/}{\texttt{rw}}), rather than with \texttt{ring} (which this book does not import from Mathlib).
